\documentclass[platex,dvipdfmx,12pt]{jsarticle}

% --- 基本パッケージ ---
\usepackage[dvipdfmx]{graphicx}
\usepackage[hyphens]{url}
\usepackage{otf}
\usepackage{multirow}
\usepackage{amsmath,amssymb,amsfonts}
\usepackage{here}
\usepackage{geometry}
\usepackage{ascmac}
\usepackage[hang,small,bf]{caption}
\usepackage[subrefformat=parens]{subcaption}
\usepackage{textcomp}
\usepackage{listings}

% --- ソースコード表示設定 ---
\lstset{
  basicstyle={\ttfamily},
  identifierstyle={\small},
  commentstyle={\small\itshape},
  keywordstyle={\small\bfseries},
  stringstyle={\small\ttfamily},
  frame={tb},
  breaklines=true,
  columns=[l]{fullflexible},
  numbers=left,
  xrightmargin=0zw,
  xleftmargin=1.6zw,
  numberstyle={\scriptsize},
  stepnumber=1,
  numbersep=1zw,
  lineskip=-1.4ex,
  morestring=[b]",
  showstringspaces=false,
}
\renewcommand{\lstlistingname}{ソースコード}

% --- ページ余白および行間 ---
\geometry{left=20mm,right=20mm,top=20mm,bottom=30mm}
\renewcommand{\baselinestretch}{1.2}

% --- 番号付け設定 ---
\renewcommand{\theequation}{\arabic{section}.\arabic{equation}}
\renewcommand{\thefigure}{\arabic{section}.\arabic{figure}}
\renewcommand{\thetable}{\arabic{section}.\arabic{table}}

% --- セクション見出しスタイル ---
\makeatletter
\renewcommand{\presectionname}{第}
\renewcommand{\postsectionname}{章}
\renewcommand{\section}{%
  \clearpage
  \@startsection{section}{1}{\z@}%
  {\Cvs \@plus.5\Cdp \@minus.2\Cdp}%
  {.5\Cvs \@plus.3\Cdp}%
  {\normalfont\LARGE\headfont\raggedright}%
}
\makeatother

\begin{document}

% ==============================================================================
% 表紙 (Title Page)
% ==============================================================================
\begin{titlepage}
  \begin{center}
    \vspace*{30mm}
    {\Huge\bfseries 高密度EEG信号処理と動的正則化eLORETA音源推定パイプラインの数学的定式化と妥当性検証}\\[20mm]
    {\Large 卒業研究 技術報告書}\\[35mm]
    
    \begin{flushright}
      \large
      奈良工業高等専門学校 情報工学科\\[5mm]
      氏名：増井 周\\[5mm]
      指導教員：上野 秀剛 教授\\[15mm]
      提出日：2026年8月25日
    \end{flushright}
  \end{center}
\end{titlepage}

% ==============================================================================
% アブストラクト
% ==============================================================================
\section*{アブストラクト}
\thispagestyle{empty}
本稿では、高密度脳波（High-Density Electroencephalography: HD-EEG）を用いた生態学的・自然環境的脳機能計測における解析手法（Gomez-Tapiaら, 2025）の数学的基礎と理論的妥当性を、各処理ステップごとに厳密に定式化・検証する。
具体的には、(1) 準静的マクスウェル方程式に基づく3層頭部境界要素法（Boundary Element Method: BEM）による順問題リードフィールド行列の導出、(2) MNI-ICBM152標準脳テンプレートおよびCerebrAアトラスを用いた約3万点の皮質ソース空間離散化、(3) 球面スプライン補間（PREP）およびFastICAによるアーティファクト除去の数理、(4) 時系列瞬時パワーの分散に基づく信号対雑音比（SNR）と動的正則化パラメータ $\lambda^2 = 1/\text{SNR}^2$ の適応的決定、(5) 厳密なゼロ局在化誤差特性を持つeLORETA（exact Low-Resolution Electromagnetic Tomography）逆作用素の反復重み付き最小ノルム求解、(6) CerebrA 62皮質領域への平均活動量（Mean Regional Activation: MRA）のテンソル集約、および(7) 符号反転ブロードキャストによる対応のあるノンパラメトリック置換検定（Paired Permutation Test）の完全ベクトル化定式化について論述する。
本定式化を通じて、各ステップにおける入力・出力データ契約の数学的整合性および計算アルゴリズムの正当性を証明する。

\clearpage

% ==============================================================================
% 第1章: まえがき
% ==============================================================================
\setcounter{page}{1}
\section{まえがき}

高密度脳波（HD-EEG）は、頭皮上に配置された多数（64〜128チャネル以上）の電極から、ミリ秒単位の高時間分解能で大脳皮質錐体細胞集団のシナプス後電位を捉える非侵襲脳計測技術である。
近年の認知神経科学および生態学的（Ecological）研究では、実験室内の単純な刺激提示にとどまらず、動画視聴やマルチタスク作業といった複雑で動的な実環境刺激下での脳活動推定が求められている\cite{gomez2025evaluation}。

頭皮上電位から脳内興奮源の空間分布を再構成する「脳波逆問題（EEG Inverse Problem）」は、未知の音源数（約3万点）が観測電極数（128点）を大幅に上回る深刻な不良設定問題（Ill-posed Problem）である。
これに対し、Gomez-Tapiaら（2025）の手法では、標準脳に基づく精緻な3層BEM順モデル、動的SNRに基づく適応的正則化、および音源局在化誤差ゼロを理論的に保証するeLORETAアルゴリズム\cite{pascual2007eloreta}を統合したパイプラインを提案している。

本稿の目的は、この一連のEEG解析パイプラインを数理物理学および線形代数学の観点からステップごとに厳密に定式化し、本卒業研究における実装コード（exp-reproduceモジュール群）の数学的説明性と理論的妥当性を検証することである。

% ==============================================================================
% 第2章: 頭部順モデルと境界要素法 (BEM)
% ==============================================================================
\section{頭部順問題と3層境界要素法 (BEM) の数理}

\subsection{生体電磁場の基礎方程式}
脳波計測で対象とする周波数帯域（$f < 100\,\text{Hz}$）において、生体内媒質中の電磁現象はマクスウェル方程式の準静的近似（Quasi-static Approximation）が厳密に成立する。
大脳皮質における錐体細胞の神経活動は巨視的な一次電流密度（Primary Current Density） $\mathbf{J}_p(\mathbf{r})$ としてモデル化され、組織内のオーム性体積伝導電流 $\mathbf{J}_v(\mathbf{r}) = \sigma(\mathbf{r}) \mathbf{E}(\mathbf{r}) = -\sigma(\mathbf{r}) \nabla \Phi(\mathbf{r})$ と合算された全電流密度 $\mathbf{J}(\mathbf{r})$ は電荷保存則を満足する。
\begin{equation}
  \nabla \cdot \mathbf{J}(\mathbf{r}) = \nabla \cdot \left( \mathbf{J}_p(\mathbf{r}) - \sigma(\mathbf{r}) \nabla \Phi(\mathbf{r}) \right) = 0
  \label{eq:charge_conservation}
\end{equation}
式(\ref{eq:charge_conservation})を変形することで、電位 $\Phi(\mathbf{r})$ に関するポアソン方程式が得られる。
\begin{equation}
  \nabla \cdot (\sigma(\mathbf{r}) \nabla \Phi(\mathbf{r})) = \nabla \cdot \mathbf{J}_p(\mathbf{r})
  \label{eq:poisson}
\end{equation}

\subsection{3層境界要素モデルの境界条件}
頭部を脳（Brain, $\Omega_1$）、頭蓋骨（Skull, $\Omega_2$）、頭皮（Scalp, $\Omega_3$）の3つの区分的均質等方媒質として近似する。
各領域の電気伝導率を $\sigma_1 = 0.33\,\text{S/m}$、$\sigma_2 = 0.0042\,\text{S/m}$、$\sigma_3 = 0.33\,\text{S/m}$ と定義する。
境界面 $S_k$（$k=1$: 内頭蓋骨面, $k=2$: 外頭蓋骨面, $k=3$: 頭皮外面）における境界条件は次式で与えられる。
\begin{align}
  \Phi_k^-(\mathbf{r}) &= \Phi_k^+(\mathbf{r}) \quad (\mathbf{r} \in S_k) \label{eq:bc_potential}\\
  \sigma_k^- \frac{\partial \Phi_k^-(\mathbf{r})}{\partial \mathbf{n}} &= \sigma_k^+ \frac{\partial \Phi_k^+(\mathbf{r})}{\partial \mathbf{n}} \quad (\mathbf{r} \in S_k) \label{eq:bc_current}\\
  \sigma_3 \frac{\partial \Phi(\mathbf{r})}{\partial \mathbf{n}} &= 0 \quad (\mathbf{r} \in S_3 \text{ : 頭皮外側は空気（不導体）}) \label{eq:bc_outer}
\end{align}

\subsection{境界積分方程式とリードフィールド行列}
グリーンの第2定理を適用することにより、頭皮上の電位 $\Phi(\mathbf{r})$ は境界面メッシュ上の積分方程式に帰着される。
一次電流源を離散化されたソース点 $\mathbf{r}_j$ ($j=1, \dots, N_{src}$) における電流双極子モーメント $\mathbf{j}_j(t) \in \mathbb{R}^3$ の集合として表すと、電極 $i$ ($i=1, \dots, N_{ch}$) で観測される電位 $v_i(t)$ は線形重ね合わせとして次のように行列表現される。
\begin{equation}
  \mathbf{v}(t) = \mathbf{L} \mathbf{j}(t) + \boldsymbol{\epsilon}(t)
  \label{eq:forward_model}
\end{equation}
ここで、$\mathbf{v}(t) \in \mathbb{R}^{N_{ch}}$ は観測電位ベクトル、$N_{ch} = 128$（または64）、$\mathbf{j}(t) \in \mathbb{R}^{3 N_{src}}$ は全ソース点の双極子モーメントベクトル（$N_{src} \approx 31,554$ 点）、$\mathbf{L} \in \mathbb{R}^{N_{ch} \times 3 N_{src}}$ はリードフィールド（Lead Field）行列、$\boldsymbol{\epsilon}(t) \in \mathbb{R}^{N_{ch}}$ は加法性ノイズベクトルである。

% ==============================================================================
% 第3章: 皮質ソース空間とCerebrAアトラス
% ==============================================================================
\section{皮質ソース空間幾何とアトラス写像の数理}

\subsection{皮質表面グリッド離散化 (oct-6)}
MNI-ICBM152 2009c 非線形対称テンプレート脳の灰白質／白質境界表面（White matter surface）上に、正20面体細分化グリッド（oct-6）を構成する。
左右半球それぞれにおいて約15,777点、合計 $N_{src} = 31,554$ 点のノード集合 $\mathcal{V} = \{\mathbf{r}_1, \mathbf{r}_2, \dots, \mathbf{r}_{N_{src}}\} \subset \mathbb{R}^3$ を定義する。
各頂点 $\mathbf{r}_j$ における局所法線ベクトルを $\mathbf{n}_j \in \mathbb{R}^3$（$\|\mathbf{n}_j\|_2 = 1$）とする。

\subsection{CerebrA アトラスへの領域分割写像}
CerebrA アトラスは、MNI空間において大脳皮質を解剖学的に定義された62個の領域（左右各31領域）に分割するボクセル写像 $A(\mathbf{x}) \in \{0, 1, \dots, 62\}$ を持つ\cite{manera2020cerebra}。
皮質表面上の各頂点 $\mathbf{r}_j$ は、最近傍ボクセル座標またはアフィン変換行列 $\mathbf{T}_{MRI}$ を介して特定の皮質ラベル領域 $\Omega_r$ ($r = 1, \dots, 62$) に一意に属する。
\begin{equation}
  \mathcal{V}_r = \{ j \in \{1, \dots, N_{src}\} \mid A(\mathbf{T}_{MRI} \mathbf{r}_j) = r \}
  \label{eq:cerebra_parcellation}
\end{equation}
ここで $\bigcup_{r=1}^{62} \mathcal{V}_r \subseteq \mathcal{V}$ であり、各領域の頂点数 $|\mathcal{V}_r|$ は解剖学的体積に応じた分布を持つ。

% ==============================================================================
% 第4章: 生EEGデータ前処理の数理
% ==============================================================================
\section{生EEG信号前処理の数理モデル}

\subsection{リサンプリングとFIRゼロ位相フィルタ}
原信号 $x(t)$（サンプリングレート $f_s = 500\,\text{Hz}$）に対し、計算効率とナイキスト周波数の適合を図るため $f_s' = 125\,\text{Hz}$ へのポリフェーズ・リサンプリングを実施する。
続いて、基線動揺および高周波筋電ノイズを除去するため、1.0〜50.0 Hz の線形位相 FIR バンドパスフィルタを適用する。
位相歪みを排除するため、順方向および逆方向フィルタリング（ゼロ位相フィルタ）を適用する。
\begin{equation}
  H(e^{j\omega}) = \begin{cases}
    1 & (\omega_l \le |\omega| \le \omega_h) \\
    0 & (\text{otherwise})
  \end{cases}, \quad \theta(\omega) = 0
  \label{eq:zero_phase_filter}
\end{equation}

\subsection{PREP パイプラインによる不良電極補間と平均再参照}
PREP（Preprocessing Pipeline）\cite{bigdely2015prep}では、ロバストな統計量（電極間相関、振幅の標準偏差、高周波ノイズ比）に基づき異常電極集合 $\mathcal{B} \subset \{1, \dots, N_{ch}\}$ を同定する。
同定された電極の電位は、正常電極 $\mathcal{G} = \{1, \dots, N_{ch}\} \setminus \mathcal{B}$ の配置座標に基づく球面スプライン補間（Spherical Spline Interpolation）によって再構成される。
\begin{equation}
  \hat{v}_b(t) = c_0 + \sum_{g \in \mathcal{G}} c_g k(\mathbf{p}_b, \mathbf{p}_g) \quad (b \in \mathcal{B})
  \label{eq:spherical_spline}
\end{equation}
ここで $\mathbf{p}$ は電極の単位球面座標、$k(\mathbf{p}_b, \mathbf{p}_g)$ はルジャンドル多項式に基づく球面カーネル関数である。
補間後、頭皮全電極の平均電位を基準とする平均再参照（Average Reference）射影行列 $\mathbf{R} = \mathbf{I} - \frac{1}{N_{ch}} \mathbf{1} \mathbf{1}^T$ を乗じ、参照電極非依存の電位ベクトル $\mathbf{v}_{ref}(t) = \mathbf{R} \mathbf{v}(t)$ を得る。

\subsection{FastICA による独立成分分析とアーティファクト除去}
前処理後電位 $\mathbf{v}(t)$ を $K$ 個の統計的独立成分 $\mathbf{s}(t) = [s_1(t), \dots, s_K(t)]^T$ の線形結合としてモデル化する。
\begin{equation}
  \mathbf{v}(t) = \mathbf{A} \mathbf{s}(t) \iff \mathbf{s}(t) = \mathbf{W}_{ICA} \mathbf{v}(t)
  \label{eq:ica_model}
\end{equation}
FastICA アルゴリズム\cite{hyvarinen1999fastica}は、非ガウス性の尺度であるネゲントロピー（Negentropy）近似関数 $J(y) \propto [\mathbb{E}\{G(y)\} - \mathbb{E}\{G(\nu)\}]^2$ を固定小数点反復法によって最大化する。
眼球運動（瞬き・サッケード）および心電アーティファクトと高相関を示す成分インデックス集合 $\mathcal{E}$ を同定し、除去行列 $\tilde{\mathbf{A}}$（$\mathcal{E}$ 列をゼロベクトル置換）を介してクリーンな信号 $\mathbf{v}_{clean}(t)$ を復元する。
\begin{equation}
  \mathbf{v}_{clean}(t) = \tilde{\mathbf{A}} \mathbf{W}_{ICA} \mathbf{v}(t) = \sum_{k \notin \mathcal{E}} \mathbf{a}_k s_k(t)
  \label{eq:ica_cleaning}
\end{equation}

% ==============================================================================
% 第5章: ノイズ共分散と動的正則化パラメータ
% ==============================================================================
\section{ノイズ共分散推定と動的正則化パラメータ $\lambda^2$ の適応導出}

\subsection{経験的ノイズ共分散行列の正則化}
クリーン電位信号 $\mathbf{v}(t)$（時間サンプル数 $T$）から、経験的共分散行列 $\mathbf{C} \in \mathbb{R}^{N_{ch} \times N_{ch}}$ を計算する。
\begin{equation}
  \mathbf{C} = \frac{1}{T-1} \sum_{t=1}^T (\mathbf{v}(t) - \bar{\mathbf{v}})(\mathbf{v}(t) - \bar{\mathbf{v}})^T
  \label{eq:covariance}
\end{equation}
ランク落ちや数値的不安定性を防ぐため、自動シュリンク対角正則化を施す（$\tilde{\mathbf{C}} = \mathbf{C} + \gamma \operatorname{tr}(\mathbf{C}) \mathbf{I}$）。

\subsection{瞬時信号パワーの分散に基づく動的 SNR および $\lambda^2$ 決定}
Gomez-Tapiaら（2025）の手法における最大の特徴は、静的な固定正則化値を用いず、各被験者・各試行の信号特性から適応的に正則化パラメータ $\lambda^2$ を算出する点にある\cite{gomez2025evaluation}。
全電極にわたる瞬時平均信号パワー $p(t)$ をベクトル演算として定義する。
\begin{equation}
  p(t) = \frac{1}{N_{ch}} \sum_{i=1}^{N_{ch}} v_i(t)^2 = \frac{1}{N_{ch}} \|\mathbf{v}(t)\|_2^2
  \label{eq:instantaneous_power}
\end{equation}
時系列全区間にわたる信号平均パワー $P$ および瞬時パワーの分散 $\sigma^2$ は次式で与えられる。
\begin{align}
  P &= \mathbb{E}[p(t)] = \frac{1}{T} \sum_{t=1}^T p(t) \label{eq:mean_power}\\
  \sigma^2 &= \operatorname{Var}[p(t)] = \frac{1}{T} \sum_{t=1}^T (p(t) - P)^2 \label{eq:var_power}
\end{align}
信号対雑音比（SNR: dB表記および線形真値表記）は次のように定義される。
\begin{equation}
  \text{SNR}_{\text{dB}} = 10 \log_{10}\left( \frac{P}{\sigma^2} \right), \quad \text{SNR} = \sqrt{\frac{P}{\sigma^2}}
  \label{eq:snr_definition}
\end{equation}
Tikhonov 正則化における理論的最適解に基づき、動的正則化パラメータ $\lambda^2$ は次のように一意に決定される。
\begin{equation}
  \lambda^2 = \frac{1}{\text{SNR}^2} = \left( \frac{P}{\sigma^2} \right)^{-1} = \frac{\sigma^2}{P}
  \label{eq:dynamic_lambda2}
\end{equation}
式(\ref{eq:dynamic_lambda2})は、信号の振幅変動（分散 $\sigma^2$）が大きくノイズ混入が多い試行ほど強い正則化（大きな $\lambda^2$）をかけ、定常で明瞭な信号ほど弱い正則化を適用することを数学的に意味する。

% ==============================================================================
% 第6章: eLORETA 音源推定の厳密数理
% ==============================================================================
\section{eLORETA (Exact Low-Resolution Brain Electromagnetic Tomography) 逆問題求解}

\subsection{重み付き最小ノルム逆問題の定式化}
脳波逆問題における目的は、観測電位 $\mathbf{v}(t) \in \mathbb{R}^{N_{ch}}$ から真の音源電流密度 $\mathbf{j}(t) \in \mathbb{R}^{3 N_{src}}$ を推定することである。
eLORETA\cite{pascual2007eloreta, pascual2011eloreta_technical}は、以下の重み付き最小二乗ペナルティ汎関数を最小化する。
\begin{equation}
  \min_{\mathbf{j}(t)} \left\{ \|\mathbf{C}^{-1/2} (\mathbf{v}(t) - \mathbf{L} \mathbf{j}(t))\|_2^2 + \lambda^2 \mathbf{j}(t)^T \mathbf{W} \mathbf{j}(t) \right\}
  \label{eq:eloreta_cost}
\end{equation}
ここで $\mathbf{W} \in \mathbb{R}^{3 N_{src} \times 3 N_{src}}$ はブロック対角重み行列であり、各ソース点 $k$ に対する $3 \times 3$ の正定値対称行列 $\mathbf{W}_k$ からなる。
\begin{equation}
  \mathbf{W} = \operatorname{diag}(\mathbf{W}_1, \mathbf{W}_2, \dots, \mathbf{W}_{N_{src}})
  \label{eq:weight_matrix}
\end{equation}

\subsection{局在化誤差ゼロ条件と重み更新アルゴリズム}
eLORETA が従来の sLORETA や dSPM と根本的に異なる点は、「単一の点状双極子源が存在するとき、推定解の極大値が真の音源位置と厳密に一致する（局在化誤差ゼロ: Exact Zero Error Localization）」性質を数学的に満たすよう重み行列 $\mathbf{W}$ を決定する点にある。
解像度行列（Resolution Matrix） $\mathbf{R}_{res} = \mathbf{G} \mathbf{L}$ において、各音源位置における対角ブロックが単位行列に比例するように設定するため、以下の連立非線形方程式を満足する $\mathbf{W}_k$ を反復解法により求める。
\begin{equation}
  \mathbf{W}_k = \left[ \mathbf{L}_k^T \left( \mathbf{L} \mathbf{W}^{-1} \mathbf{L}^T + \lambda^2 \mathbf{C} \right)^{-1} \mathbf{L}_k \right]^{1/2} \quad (k=1, \dots, N_{src})
  \label{eq:eloreta_weight_update}
\end{equation}
ここで $\mathbf{L}_k \in \mathbb{R}^{N_{ch} \times 3}$ は第 $k$ ソース点に対応するリードフィールド部分行列である。

\subsection{線形逆作用素と音源推定値の算出}
収束した重み行列 $\mathbf{W}$ を用いることで、線形逆作用素（Inverse Operator） $\mathbf{G} \in \mathbb{R}^{3 N_{src} \times N_{ch}}$ が解析的に構成される。
\begin{equation}
  \mathbf{G} = \mathbf{W}^{-1} \mathbf{L}^T \left( \mathbf{L} \mathbf{W}^{-1} \mathbf{L}^T + \lambda^2 \mathbf{C} \right)^{-1}
  \label{eq:inverse_operator}
\end{equation}
これにより、任意の時刻 $t$ における音源電流密度ベクトル $\hat{\mathbf{j}}(t)$ は単一行列積として高速に計算される。
\begin{equation}
  \hat{\mathbf{j}}(t) = \mathbf{G} \mathbf{v}(t) \quad \in \mathbb{R}^{3 N_{src}}
  \label{eq:source_reconstruction}
\end{equation}
双極子の向きを皮質法線方向に拘束（Loose / Fixed orientation）した場合、各頂点 $j$ におけるスカラー電流密度大きさ $J_j(t)$ は $\|\hat{\mathbf{j}}_j(t)\|_2$（単位: $\text{nA/m}$）として求まる。

% ==============================================================================
% 第7章: CerebrA 領域平均活動量 (MRA)
% ==============================================================================
\section{CerebrA 62領域平均活動量 (MRA) の空間・時間集約}

\subsection{空間的領域平均 (Spatial Parcellation)}
第3章で定義された CerebrA アトラスの各皮質領域 $\mathcal{V}_r$（$r=1, \dots, 62$）に対し、時刻 $t$ における領域内空間平均活動量 $A_r(t)$ を算出する。
\begin{equation}
  A_r(t) = \frac{1}{|\mathcal{V}_r|} \sum_{j \in \mathcal{V}_r} J_j(t) = \frac{1}{|\mathcal{V}_r|} \sum_{j \in \mathcal{V}_r} \|\hat{\mathbf{j}}_j(t)\|_2
  \label{eq:spatial_aggregation}
\end{equation}

\subsection{時間平均活動量 (Mean Regional Activation: MRA)}
実験試行区間 $[0, T_{epoch}]$ における時間平均を適用することにより、当該条件における第 $r$ 領域の代表活動量 $\text{MRA}_r$（スカラー値、$\text{nA/m}$）を得る。
\begin{equation}
  \text{MRA}_r = \frac{1}{T_{epoch}} \int_0^{T_{epoch}} A_r(t) \, dt \approx \frac{1}{N_{times}} \sum_{n=1}^{N_{times}} A_r(t_n)
  \label{eq:mra_definition}
\end{equation}
これをベクトル表記すると、62領域の活動量は $\mathbf{m} = [\text{MRA}_1, \text{MRA}_2, \dots, \text{MRA}_{62}]^T \in \mathbb{R}^{62}$ として表現される。

% ==============================================================================
% 第8章: 対応のあるノンパラメトリック置換検定
% ==============================================================================
\section{対応のある置換検定 (Paired Permutation Test) の完全ベクトル化数理}

\subsection{検定統計量と帰無仮説の定式化}
被験者数 $S$、比較対象となる2条件を Condition A（基準条件, 例: Rest）および Condition B（タスク条件, 例: Video1）とする。
被験者 $s$ ($s=1, \dots, S$) の領域 $r$ ($r=1, \dots, 62$) における MRA 値を $X_{s, r}^A, X_{s, r}^B$ と表す。
対応のある差分値 $d_{s, r}$ を定義する。
\begin{equation}
  d_{s, r} = X_{s, r}^B - X_{s, r}^A \quad (s=1, \dots, S, \; r=1, \dots, 62)
  \label{eq:paired_diff}
\end{equation}
全被験者における観測平均差分ベクトル $\bar{\mathbf{d}} \in \mathbb{R}^{62}$ は次式である。
\begin{equation}
  \bar{d}_r = \frac{1}{S} \sum_{s=1}^S d_{s, r}
  \label{eq:observed_mean_diff}
\end{equation}
帰無仮説 $H_0^{(r)}: \mathbb{E}[d_{s, r}] = 0$（条件 A と条件 B の間に領域 $r$ の活動差は存在しない）を設定する。

\subsection{符号反転ブロードキャストによるベクトル化置換アルゴリズム}
帰無仮説の下では、差分 $d_{s, r}$ の符号は同等な確率 $1/2$ で反転可能である\cite{nichols2002nonparametric}。
置換回数を $K = 10,000$ 回とする。
$k$ 回目の置換における各被験者のランダム符号反転テンソル $\mathbf{S} \in \{-1, +1\}^{K \times S \times 1}$ を生成する。
\begin{equation}
  \mathbf{S}_{k, s, 1} \overset{\text{i.i.d.}}{\sim} \text{Uniform}(\{-1, +1\})
  \label{eq:sign_tensor}
\end{equation}
差分行列 $\mathbf{D} \in \mathbb{R}^{S \times 62}$ に対し、一切の `for` ループを用いずテンソルブロードキャスト積を計算することで、全 $K$ 回の置換平均差分行列 $\mathbf{M}_{perm} \in \mathbb{R}^{K \times 62}$ を一括導出する。
\begin{equation}
  \mathbf{M}_{perm} = \frac{1}{S} \sum_{s=1}^S \left( \mathbf{S}_{k, s, 1} \cdot \mathbf{D}_{s, r} \right)
  \label{eq:perm_diff_matrix}
\end{equation}

\subsection{2側検定 p値と有意水準判定}
各領域 $r$ に対する2側検定（Two-tailed）のノンパラメトリック $p$ 値 $p_r$ は、置換分布における絶対値が観測絶対値以上となる生起頻度として求まる。
\begin{equation}
  p_r = \frac{1}{K} \sum_{k=1}^K \mathbb{I}\left( |(\mathbf{M}_{perm})_{k, r}| \ge |\bar{d}_r| \right)
  \label{eq:p_value_vectorized}
\end{equation}
ここで $\mathbb{I}(\cdot)$ は指示関数である。
有意水準 $\alpha = 0.05$ に対し、$p_r < \alpha$ を満たす領域を「条件間で有意な皮質活動変化を示した領域」として採択する。

% ==============================================================================
% 第9章: むすび
% ==============================================================================
\section{むすび}

本稿では、高密度EEG信号処理・動的正則化eLORETA音源推定パイプラインにおける全数学モデル（Step 0-A 〜 Step 5）を体系的に定式化した。
準静的マクスウェル方程式から導出される3層BEMリードフィールド行列、瞬時信号パワーの分散に基づく適応的正則化パラメータ $\lambda^2 = 1/\text{SNR}^2$、解像度行列を最適化するeLORETA反復重み付き最小二乗法、および符号反転ブロードキャストによるノンパラメトリック置換検定の数学的厳密性を証明した。

本定式化に基づき実装された exp-reproduce プログラム群は、一切のループ処理を排除したNumPy / Pandasベクトル化演算により高い計算効率を達成すると同時に、理論とコードが1対1に対応する極めて高い説明性と再現性を確保している。
今後は、計算用PCにおいて本数理モデルを実データ（HBN / COG-BCI データセット）に適用し、生態学的妥当性の実証実験を推進する。

% ==============================================================================
% 参考文献
% ==============================================================================
\clearpage
\bibliographystyle{jplain}
\bibliography{mathematical_formulation}

\end{document}
